
\subsubsection{6.1 Unexpected Detection Rate}

The 99.6\% mypy detection rate far exceeds the predicted 30-40\% range. Several factors may contribute:

\begin{enumerate}
\item \textbf{Task selection bias:} Dual-sensitive classification may have selected tasks where CodeLlama-7B systematically generates type annotation issues
\item \textbf{Base model characteristics:} The non-instruction-tuned model may lack type safety training, producing pervasive static analysis errors
\item \textbf{Mypy --strict comprehensiveness:} Strict mode may catch broader issue classes than anticipated (unreachable code, missing return statements, beyond pure type errors)
\end{enumerate}

The high detection rate validates the computational efficiency justification for cascade routing: nearly all iterations can skip expensive execution when static errors exist.

\subsubsection{6.2 Attention Economy Unverified}

The H-M2 implementation failure leaves the attention economy hypothesis untested. Without data comparing iteration counts between sequential and aggregation presentation, we cannot claim that cascade routing reduces iterations to solution.

The validated contribution is limited to computational efficiency (early error detection with zero execution cost) rather than cognitive efficiency (iteration reduction via single-source focus).

\subsubsection{6.3 Scope and Generalization}

Results apply specifically to:
\begin{itemize}
\item Dual-sensitive tasks (21.3\% of HumanEval, N=35)
\item CodeLlama-7B base model (not instruction-tuned variants)
\item Python code with mypy --strict static analysis
\item HumanEval benchmark programming problems
\end{itemize}

Generalization to larger models, instruction-tuned variants, other programming languages, or non-dual-sensitive tasks requires separate validation. The 99.6\% detection rate may not hold for models with stronger type annotation capabilities.

\subsubsection{6.4 Mock vs. Real Validation}

The token efficiency finding (H-M3) is based on mock simulation. While the routing logic was verified correct and all code paths tested, actual token counts with CodeLlama-7B inference may differ from simulated values. The 0.733 ratio should be interpreted as plausible given the 99.6\% mypy detection rate (which enables extensive execution skipping), but not empirically confirmed.

\subsubsection{6.5 Comparison to Existing Work}

The study confirms that static analysis can provide substantial early error detection in LLM code generation feedback loops, extending prior work that demonstrated static analysis reduces error rates in post-generation validation. The 99.6\% detection rate suggests static analysis may be nearly universal as a first-pass filter for certain model-task combinations.

The cascade vs. aggregation comparison addresses a gap in prior literature, which has not systematically tested feedback routing policies with controlled ablation methodology.
